\section{A bank funding decision with common shocks and learned controls}
\label{s6:funding}
Consider a treasury desk changing the offered deposit rate for one month by
ten basis points. The target is the response of a fixed branch portfolio's
balance stock, followed by the change in funding contribution over seven
months. This is neither a sustained repricing policy nor the response to an
unlabelled reduced-form innovation. Future prices follow their natural rule.
The customer and competitor mechanisms are held fixed; equilibrium repricing
by competitors would require a different model.

For branch $i$, write $B_{it}$ in millions of currency units, $A_t$ in units
of ten basis points, and $W_t$ for an observed common funding driver. A
representative construction is
\begin{align}
 W_t&=\phi W_{t-1}+\sqrt{1-\phi^2}\,\zeta_t,\nonumber\\
 A_t&=a(W_t)+\nu_t,\qquad a(w)=.5w+.2(w^2-1),\nonumber\\
 B_{it}&=\rho B_{i,t-1}+\beta_i A_t+g(W_t)+\eta_t+\epsilon_{it},
 \quad g(w)=1.2w+.4(w^2-1).\label{s6:funding-law}
\end{align}
All innovation sequences are independent Gaussian sequences, with
$\nu_t,\zeta_t$ standardized to variance one; $\eta_t$ is
shared by all branches. Stability margins for $|\phi|,|\rho|<1$, finite
moments and a fixed finite horizon are part of the experiment. The
simulation has 48 equally weighted branches, $\rho=.82$, $\phi=.65$,
$\bar\beta=2$, common noise standard deviation 2 and individual noise
standard deviation 3. Heterogeneous $\beta_i$ range from 1.2 to 2.8.
Adding the baseline stock of 100 per branch changes accounting but not
response estimation.
The simulation starts its centered states at zero and discards 500 dates;
it approximates the stationary law rather than sampling its initial law
exactly. At the declared stability margins the remaining initialization
effect decays geometrically. This construction cannot justify using the
same burn-in for a process arbitrarily close to a unit root.

Let $\bar B_t=\sum_iw_iB_{it}$ for fixed weights summing to one. Because
the intervened price does not change $W$ or future assignment, recursive
substitution gives
\[
 \theta_h=\E[\bar B_{t+h}(A_t+1)-\bar B_{t+h}(A_t)]
       =\bar\beta\rho^h,\qquad h=0,\ldots,H.
\]
This is an additive shift of a continuous action with support on the real
line in the reference law. A bounded bank price grid needs its own feasible
shift and support restriction; Gaussian prices are not a literal product
pricing model. Identification uses the independence of $\nu_t$ from future
structural disturbances conditional on
$X_t=(\bar B_{t-1},W_t)$. That exclusion is stronger than a good prediction
of prices. It fails if officers privately anticipate future withdrawals.
An external instrument is then a substantive alternative only if its own
exclusion and normalization can be defended.

\subsection{The estimator, learning remainder and calendar experiment}
Define $m_h(X)=\E[\bar B_{t+h}\mid X]$ and $e(X)=\E[A_t\mid X]$.
For $v_t=A_t-e(X_t)$, the structural law implies
\[
 \E[v_t\{\bar B_{t+h}-m_h(X_t)-\theta_hv_t\}]=0,
 \qquad q=\E v_t^2=1.
\]
The finite basis $(1,\bar B_{t-1},W_t,W_t^2-1)$ contains these nuisance
functions under Equation~\eqref{s6:funding-law}: Gaussian AR conditional
means of future linear and quadratic driver terms stay in this basis.
The partialling-out construction is a selected specialization of the
orthogonal approach in \citet{dml2018}; the time-series extension in
\citet{ciganovic2026} requires its additional stability and sampling
conditions, not merely chronological input rows.

An independent pilot realization chooses a ridge penalty from
$\{0,1,10,100\}$ by a chronological half-pilot validation loss, then fits
all nuisance functions on the pilot. The family, candidate penalties and
normalization are fixed before that pilot. The evaluation sample never
selects a penalty, horizon, regime or bandwidth. Independent realization
means an independent statistical population under the same law; an adjacent
historical block is not literally independent under this model. In a bank,
using a separated block instead requires a gap/mixing argument and its
finite-sample sensitivity. Pooling another country additionally requires
transport of the nuisance law. The reference deliberately makes this
information cost visible.

With frozen pilot functions $\hat e,\hat m_h$, set
$\hat v=A-\hat e$ and estimate
\[
 \hat\theta_h=
 \frac{\mathbb P_T[\hat v(\bar B_{t+h}-\hat m_h)]}
      {\mathbb P_T\hat v^2}.
\]
Write $d_e=\hat e-e$, $d_h=\hat m_h-m_h$ and
$u_h=\bar B_{t+h}-m_h-\theta_hv$. Conditional on the pilot, the
expectation of the estimating score at $\theta_h$ is exactly
\begin{equation}
 \E[d_e(X)d_h(X)]-\theta_h\E[d_e(X)^2].
 \label{s6:pilot-remainder}
\end{equation}
To see this, expand
$(v-d_e)(u_h-d_h+\theta_hd_e)$; conditional means of $v$ and $u_h$
vanish. Thus independence alone does not remove learning bias. Sufficient
conditions are nuisance consistency and
$\|d_e\|_2(\|d_h\|_2+|\theta_h|\|d_e\|_2)=o_p(T^{-1/2})$,
with corresponding moment/stochastic-equicontinuity bounds. The fixed
correct basis, independent pilot size proportional to $T$ and bounded
penalty candidates give root-pilot coefficient error; the remainder is
$O_p(T^{-1})$. An arbitrary forest or representation learner would need
a new rate argument.

Under the stable Gaussian-driven polynomial law, covariances and the
required mixing/moment sums decay geometrically. The vector influence is
\[
 \psi_{th}=q^{-1}v_tu_{th},\qquad
 \sqrt T(\hat\theta-\theta)
 =T^{-1/2}\sum_t\psi_t+o_p(1).
\]
A mixing central limit theorem therefore gives covariance
$\Omega=\sum_{k\in\mathbb Z}\E[\psi_t\psi_{t-k}']$.
The evaluation residuals replace both $\theta$ and the pilot nuisances in
$\hat\psi$. Bartlett covariance with
$L=\lfloor2T^{1/5}\rfloor$ uses the denominator $T$ at every lag and
estimates $\Omega/T$ \citep{newey1986}. For the fixed basis and moments
here, coefficient replacement contributes at most
$O_p(LT^{-1/2})$ to the lag-sum bound, and the taper and tail biases vanish.
The exponent is strictly below the general source's $1/4$ bound.
This argument is pointwise in the stability margin and fixed $H$; it gives
no uniform near-unit-root guarantee. Overlapping horizons share future
outcomes, and their covariance must remain in $\hat\Omega$.

Increasing branches removes part of $\epsilon_{it}$, but not $\eta_t$ or
the common action innovation. Calendar origins are the information count.
Dividing the calendar variance by the number of branches is an adverse
comparison, not another interpretation of cluster robustness. Bonferroni
normal intervals give a conservative joint response band; a prespecified
financial contrast uses its own scalar normal interval. Simulated coverage
checks the chosen regime, not the asymptotic assumptions.

\subsection{Funding accounting and method decisions}
Let $N$ be the branch count, $c=10^6$ the currency conversion, $m=.002$
the monthly net alternative-funding margin and $d=.99$ the discount factor.
The first-month additional rate is $r=.001/12$. Relative to no shift,
\begin{equation}
 V=Nc\left\{\sum_{h=0}^H d^h m\theta_h
                 -r(B_0+\theta_0)\right\},\qquad
 \widehat{\Var}(\hat V)=k'\hat\Omega k/T,
 \label{s6:funding-account}
\end{equation}
where $k_h=Nc(d^hm-r\mathbf1\{h=0\})$ and $B_0=100$ is fixed.
The higher rate also applies to induced first-month balances. Omitting that
term overstates value. There is no cost saving from relabelling a balance
stock as a transaction flow. If baseline exposure, margin or discounting is
estimated, their influence and covariance must be added; the displayed
interval conditions on these economic inputs.

The primary procedure is the pilot-trained calendar score because its
functional, learning error and financial covariance can all be derived for
this information regime. Direct residual regression with the same frozen
pilot and Bartlett rule is a computational comparator: it solves precisely
the same score. An outcome regression fitted jointly on evaluation data can
be more efficient in the correctly specified finite basis, but is a
different fitted experiment. External-IV projections become the primary
alternative when the price innovation is privately confounded and a valid
instrument is available; their weak normalization needs the set-valued
inference already derived. Lag-augmented innovation projections do not by
themselves solve that IV problem.

A short bank panel with endogenous lagged balances calls for the separate
initial-law and lag-moment analysis, not multiplication of the calendar
sample size. A mixed-frequency analysis can retain daily prices and monthly
stock measurements, but must specify whether the monthly stock sees the
last daily action, its path or an integral. A MIDAS distributed-lag
restriction reduces coefficients only after that observation equation is
fixed. Synthetic controls address a sustained unit intervention under donor
restrictions; they do not automatically recover this one-month innovation
response. Each is a substantive alternative when its observation and
assignment premises match the bank's question.

If evaluation outcomes suggest a break, choosing the favorable subperiod and
reusing this interval is invalid. The result must become exploratory, or
the bank must define a new regime-specific target and obtain independent
evaluation information. The public operation binds the pilot, calendar,
bandwidth, observation mapping, price normalization and financial weights
to its inference identity. This is a necessary consequence of the derivation,
not merely a software cache preference.
